% !TeX root = ../papers/nc-p3.tex
\section{九次元半周期成分の四直線枠・三重みと二次の中心切断}
\label{sec:half-period-frame}
\begin{theorem}\label{thm:half-period-frame}
九次元半周期成分の稠密開部分には次の表示がある。四空間 $E$ に非退化対称形式 $q$ と
その二つの直線族の順序を取り、$E^*$ の順序付き四直線枠の代表を $u_i$ とする。
$q=\sum_{i=1}^4 u_i\otimes v_i$ で $v_i$ を定め、$p_i=u_i\otimes v_i$ と置く。
$\sum_i\omega_i=1$ の非零重みを取り、三平面
\[
 I=\{\textstyle\sum_i c_i p_i:\sum_i\omega_i c_i=0\}
\]
を作る。$L_+,L_-\subset\bigwedge^2 E^*$ を、$q$ の各直線族のPlücker座標上で
消える交代形式の三空間とする。二周期窓の関係は
\[
 R_0=I\oplus L_+,\qquad R_1^{\mathrm{op}}=I\oplus L_- .
\]
一般の窓から $[q]$、直線族、指定なしの四組 $(p_i,\omega_i)$ を共通射影基底変更まで回収できる。
二周期での $q$ の極化は中央の二次切断である。
八つの固定点は三つの二次式による $\mathbb P^3$ の被約完全交叉である。
五次元半周期成分の一般点には、可逆な二位置移動と交換する非零の二次切断は存在しない。
\end{theorem}
\begin{proof}[点加群と四空間の同定]
定理\ref{thm:half-period-components}の九次元枝の一般点で位数2の持上げを選んで
二周期化する。点の座標に作用する対称な第一行列を $M$ とし、
$N_{ij,kl}=M_{kj,il}$ を第二の逆順行列とする。第一辺の二十点 $(a,b)$ に
次の辺をたどって偶位置の点 $c$ を対応させる。半周期で同定すれば $a\mapsto c$ は対合となる。
実際四位置をたどった同定は有限四辺の四移動の恒等性による。
\path{computations/verify_half_period_frame.py} は53のWalsh点で
位数2持上げに固定八点・二巡回六組、三つの位数4持上げに二巡回二組・四巡回四組を得る。
六関係の有限エタールな交点を番号付けすれば置換は局所定数であり、
これらの巡回型は対応する標数零の各成分の一般点へ移る。
八つの固定点の両射影は各々射影枠を含む。五次元側では二つの持上げは互いに逆で、双方とも固定点がない。
移動と交換する切断の各点での二次評価が非零なら、三文字での交換式に
全ての一文字を代入して、二位置先の点と移動で運んだ始点とが一致する。
実際、両端の一次汎関数は非零で比例する。この比較は交換時の位置別スカラーを許しても成立する。
二次評価は六関係による十次元商上で単射である（53の評価階数10は開条件）。
従って五次元側でそのような切断は零である。

九次元側で $H=M-N$ を $\bigwedge^2 E\to(\bigwedge^2 F)^*$ の写像と読む。
二つの偶・奇点空間をまだ $E,F$ と区別している。$H$ は両添字対で交代的であり、
53の点でこの写像は同型である。固定点 $(a,b)$ では
$M(a\otimes b)=N(a\otimes b)=0$ なので、$H(a\wedge E)$ は
$b\wedge F$ の零化空間である。体積形式で双対を同定するとこれは $b\wedge F$ 自身となる。
従って $H$ によるPlücker二次形式の引戻しは八つの $a\wedge E$ の上で消える。
既に射影枠の五点でこの形式は定数倍に一意である。実際、最初の四点を座標点にすると
相補な二添字の積三つだけが残り、第五点を $(1,1,1,1)$ にして
$(1,1,1,1)\wedge z$ を代入すると係数比はPlücker式のものになる。
よって $H$ はKlein二次超曲面の射影同型で一方の平面族を保つ。
その平面族のパラメータ $\mathbb P(E)$ と $\mathbb P(F)$ の同型は射影線型で、
各固定 $a$ を対応する $b$ に送る。これで両空間を同定すると、
各 $a\wedge E$ を保つ写像は二平面の交線も保ち、全分解可能線を固定するのでスカラーである。
元の $H_{ij,kl}$ は従って体積テンソルの非零倍である。

\end{proof}
\begin{proof}[八点と二つの直線族]
$S=\mathrm{Sym}^2E$, $T=\bigwedge^2 E$ として $E\otimes E=S\oplus T$ に分ける。
$H$ の核は $S$、$T$ への制限は非退化である。
$r=\operatorname{rank}(M|_S)$ とすると、対称性で
$\dim(\operatorname{im} M\cap T^*)=6-r$ である。
よって $\operatorname{im}M+\operatorname{im}N=\operatorname{im}M+T^*$ は次元 $6+r$、
両像の交叉は次元 $6-r$ で、そこに $M(S)=N(S)$ が含まれる。従って $r\le3$。
対角上の第一辺は $r$ 本以下の二次式で定まる。その全ての点は $N$ も消すから
ちょうど固定八点である。有限性から $r=3$、Bézoutにより三つの二次式の被約完全交叉である。
一般点で対称対角ブロックの階数も3である（53で $S$ 上の階数3）。
共通像 $I=M(S)$ は二つの像の交叉で、$S$ への制限が単射となる。

$S$ の三次元非退化商でSchur消去すれば
$M=X+D_+$, $N=X+D_-$ と書ける。$X$ は対称階数3で像は $I$、
$D_+,D_-$ は $T$ 上の対称階数3で、$D_+-D_-=H$ である。
$H^{-1}D_+$ と $-H^{-1}D_-$ は和が恒等で階数の和が6であり、互いに相補な自己共役射影である。
従って両像はPlücker対合で非退化な相補直交三空間となる。
各三空間が $T$ で消す三平面とKlein二次超曲面の交叉は滑らかな平面円錐曲線である。
これは $\mathbb P^3$ の滑らかな二次曲面の一方の直線族を与え、直交する他平面は反対の直線族を与える。
実際、円錐曲線上の異なる二点の極対合は非零なので対応する直線は交わらない。
三つの互いにねじれた直線を選べば一つの滑らかな二次曲面の同じ直線族に入り、
この三直線を全て横切る直線は直交平面の円錐曲線そのものである。二度適用すると元の族を得る。
この二次曲面を $q=0$、両像を $L_+,L_-$ と書く。一般点では $q$ は固定八点の全てで非零である。
これは53で構成した対角係数 $(37,15,47,1)$ の形式で確かめられる。
\end{proof}

\begin{proof}[四直線枠の回収と逆構成]
$q$ を行列で表し、$G_{ij,kl}=q_{ik}q_{jl}-q_{il}q_{jk}$ とする。
$H=h\epsilon$ を固定体積で表す。二つの直線族の順序に合わせて
$c=\pm h/\sqrt{\det q}$ の符号を選ぶと
\[
 D_+=(H+cG)/2,\qquad D_-=(-H+cG)/2 .
\]
これは $H^{-1}cG$ の $T$ 上の二固有空間が二つの直線族の三空間となることによる。
例えば $q$ を単位行列にすれば、対合は相補二添字とその並べ替え符号を交換するHodge作用である。
$\Gamma$ を $i,k$ の部分転置とする。$Y=X+(c/2)q\otimes q$ では $Y^\Gamma=Y$ となる。
実際 $G^\Gamma-G=(q\otimes q)^\Gamma-q\otimes q$ を $M^\Gamma=N$ に代入すればよい。
一般に $Y$ の階数は4である。$q$ が固定点で非零なので $q\notin I$ である。

$Y$ の像の基底となる四つの双一次式の左変数を $a$ で評価した $4\times4$ 行列を $B(a)$ とし、
$Y$ はこれらの形式の積に非退化対称係数行列 $D$ を付けて書く。
部分転置不変性は $B(a)^{\mathsf t}D B(b)=B(b)^{\mathsf t}D B(a)$ である。
$B(a_0)$ が可逆なら $C(a)=B(a_0)^{-1}B(a)$ は
$B(a_0)^{\mathsf t}D B(a_0)$ について自己共役である。同じ式から全ての $C(a)$ は可換である。
53の検算では $a_0=(1,1,1,1)$ で行列式3、
$a=(7,4,2,9)$ で $C(a)$ の固有値は $-3,-4,-15,-52$ の四単根である。
従って一般には同時固有四直線があり、非退化対合で直交する。
相手側の四形式の基底も変えると $B(a)$ は対角となり、像の基底は四つの純双一次式になる。
その空間に非退化な $q$ が入るため両側の四因子がそれぞれ基底であり、
$q=\sum p_i$ で純形式の倍率を一意に定められる。純形式の直線はこの四本だけである。
$Y=(c/2)\sum_i p_i\otimes p_i/\omega_i$ と書くと全 $\omega_i\ne0$ で、
$X=Y-(c/2)q\otimes q$ の階数3は行列式補題から $\sum_i\omega_i=1$ と同値である。
その像が所定の $I$ である。

逆に任意のこの枠と非零重みから上式で $X,D_+,D_-$ を作れば $M^\Gamma=N$ で、
$I$ の交代三空間との交叉が零なら両隣接階数は6である。
識別に使った位数2の持上げは九次元側の一般点では射影的に一意である。
八点の両射影の同定、二つの交代三空間、$q$、$I$ を順に回収し、
$I+kq$ の四純直線と $I$ の方程式から四線と重みが回収される。
したがってこの開集合では同時添字置換と直線族の順序を保つ射影同型が同型判定を与える。
$q$ を固定して四線枠に $15-3=12$、その射影直交群に6、重みに3の自由度があり、
既知の成分次元9と一致する。最後に $R_0,R_1^{\mathrm{op}}$ の共通平面から、
偶・奇双方で $p_i=\omega_i q$（奇では転置）が成り立つ。
$(u_i v_i)u_i=u_i(v_i u_i)$ とその逆により、$q$ は全 $u_i,v_i$ と二周期で交換する。
\end{proof}
